\section{The exact cubic matching theorem}\label{sec:cubic}

The article asks whether its quadratic Euler-factor matching results have an
analogue for characters of odd prime order \cite{article}. We determine the
minimum disagreement density for characters of order dividing three, for every
packet size and every rank, and classify every pair attaining that minimum.
The resulting function has two different steps between successive powers of
three. The extremal configurations also give exact erasure and corruption
thresholds for identifying a product from its local factors.

\subsection{Packets and the extremal function}

Let $X$ be a group of Dirichlet characters modulo $Q$, of exponent three and
rank $r\geq1$. Thus $X\simeq\F_3^r$. The principal character belongs to $X$.
We write the group law on $X$ additively when describing its subspaces and
cosets. A \emph{packet} $A$ is a multiset of $k\geq1$ members of $X$, and its
local polynomial at $p\nmid Q$ is
\[
 P_A(p,T)=\prod_{\chi\in A}(1-\chi(p)T).
\]
For distinct packets $A,C$ of the same size, let $\delta(A,C)$ be the natural
density, among rational primes, of the set on which $P_A(p,T)\ne P_C(p,T)$.
Write $\Delta_{r,k}$ for the minimum of these densities over all such pairs.

\begin{theorem}[Exact minimum disagreement]\label{thm:cubic-min}
The density $\delta(A,C)$ exists. Put $a=\lfloor\log_3 k\rfloor$. Then
\begin{equation}\label{eq:cubic-min}
 \Delta_{r,k}=
 \begin{cases}
  2/3^r,&k\geq3^{r-1},\\[2pt]
  2/3^{a+1},&a\leq r-2,\quad 3^a\leq k<2\cdot3^a,\\[2pt]
  4/3^{a+2},&a\leq r-2,\quad 2\cdot3^a\leq k<3^{a+1}.
 \end{cases}
\end{equation}
The theorem is unconditional.
\end{theorem}

For example, at rank at least three, the least disagreement densities for
$k=1,2,3,4,5,6,7,8,9$ are respectively
\[
 \frac23,\quad\frac49,\quad\frac29,\quad\frac29,\quad\frac29,
 \quad\frac4{27},\quad\frac4{27},\quad\frac4{27},\quad\frac2{27}.
\]
Consequently two distinct products of six factors can have identical local
polynomials at a proportion $23/27$ of the primes, and no larger proportion
is possible. With nine factors the optimal proportion is $25/27$.

For the classification, cancel common entries of $A$ and $C$, and write
\[
 c(\chi)=m_A(\chi)-m_C(\chi),\qquad
 b=\sum_{c(\chi)>0}c(\chi)=\frac12\sum_\chi\abs{c(\chi)}\leq k.
\]
Here $b$ is the size of either packet after cancellation.

\begin{theorem}[All extremizers]\label{thm:cubic-equality}
Every pair attaining \eqref{eq:cubic-min} has the following form, up to
interchanging the two packets and adding the same multiset to both.
\begin{enumerate}
\item If $a\leq r-2$ and $3^a\leq k<2\cdot3^a$, there are a subspace
  $U\leq X$ of dimension $a$, an element $\eta\in X$, and $\theta\notin U$
  such that the reduced packets are
  \[
   A=\eta+U,\qquad C=\eta+\theta+U.
  \]
  In particular, $b=3^a$ and neither reduced packet has repeated entries.
\item If $a\leq r-2$ and $2\cdot3^a\leq k<3^{a+1}$, there are a subspace
  $U\leq X$ of dimension $a$, an element $\eta\in X$, and $\alpha,\beta$
  linearly independent modulo $U$ such that
  \begin{align*}
   A&=(\eta+\alpha+U)\mathbin{\dot\cup}(\eta-\alpha+U),\\
   C&=(\eta+\beta+U)\mathbin{\dot\cup}(\eta-\beta+U).
  \end{align*}
  Here $b=2\cdot3^a$, again without repetitions.
\item If $k\geq3^{r-1}$, there is a hyperplane $U<X$ such that $c$ is
  constant on its three cosets. Its three integer values $d_0,d_1,d_2$
  satisfy
  \[
   d_0+d_1+d_2=0,\qquad(d_0,d_1,d_2)\ne(0,0,0),\qquad
   b=3^{r-1}\max_j\abs{d_j}\leq k.
  \]
  Every such coefficient function is an extremizer.
\end{enumerate}
Conversely every displayed construction attains the stated minimum.
\end{theorem}

The third case covers $r=1$ and every $k\geq1$. The theorem concerns the exact
extremal function and its complete equality classification. The other
attainable densities form a further classification problem.

\begin{corollary}[Characters of exact order three]\label{cor:cubic-nonprincipal}
Formula~\eqref{eq:cubic-min} remains valid if every packet entry must be
nonprincipal, hence of exact order three.
\end{corollary}

\begin{proof}
The lower bound already applies to this restricted family. For the lower
construction use the two cosets $\theta+U$ and $-\theta+U$, where
$\theta\notin U$; neither contains the identity. For the upper
construction the four cosets $\alpha+U$, $-\alpha+U$, $\beta+U$ and
$-\beta+U$ already avoid the identity. At saturation take the two
nonzero cosets of a hyperplane. Pad both packets by the same nonprincipal
character as necessary; one exists because $r\geq1$. This also handles
$r=1$. The equality classification is the restriction of
Theorem~\ref{thm:cubic-equality} to configurations and common fillers
that avoid the identity.
\end{proof}

\subsection{From local factors to an integral Fourier transform}

Fix $\zeta=e^{2\pi i/3}$, and let
$V=\operatorname{Hom}(X,\{1,\zeta,\zeta^2\})\simeq\F_3^r$.
For $v\in V$ write $\chi(v)$ for evaluation. Define
\begin{equation}\label{eq:cubic-H}
 H(v)=\sum_{\chi\in X}c(\chi)\chi(v).
\end{equation}
Finite Fourier inversion shows that $H\ne0$. Also
\begin{equation}\label{eq:cubic-symmetry}
 H(0)=0,\qquad H(-v)=\overline{H(v)}.
\end{equation}

\begin{lemma}[Trace detects the whole cubic local polynomial]
\label{lem:cubic-trace}
At every $v\in V$, equality of the two traces is equivalent to equality of
the two local polynomials. Moreover
\begin{equation}\label{eq:cubic-density}
 \delta(A,C)=3^{-r}\abs{\supp H}.
\end{equation}
\end{lemma}

\begin{proof}
Let
\[
 C_j(v)=\sum_{\chi(v)=\zeta^j}c(\chi)\qquad(j=0,1,2).
\]
Then $\sum_jC_j(v)=0$ and
$H(v)=C_0(v)+C_1(v)\zeta+C_2(v)\zeta^2$.
The minimal polynomial $1+Z+Z^2$ implies that this expression vanishes
precisely when $C_0(v)=C_1(v)=C_2(v)$; their zero sum then makes all three
zero. This says exactly that the multiplicities of the three possible
local roots agree.

The evaluation homomorphism $(\Z/Q\Z)^\times\to V$ is surjective. Indeed,
the annihilator of its image in $X$ consists of the character trivial on
every reduced residue class, which is the identity character. A proper
subgroup of a finite elementary abelian group has a nontrivial annihilator.
All fibers therefore have the same size. The prime number theorem in
arithmetic progressions gives \eqref{eq:cubic-density}.
\end{proof}

\subsection{The Eisenstein valuation and the degree bound}

Put $R=\Z[\zeta]$ and $\varpi=1-\zeta$. We use
\[
 \varpi^2=-3\zeta,\qquad R/(\varpi)\simeq\F_3,
 \qquad\overline\varpi=-\zeta^2\varpi.
\]
For a nonzero element of $R$, let $v_{\varpi}$ be its divisibility order at
$\varpi$. Set
\[
 t=\min_{H(v)\ne0}v_{\varpi}(H(v)),\qquad
 f(v)=\varpi^{-t}H(v)\pmod\varpi.
\]
The function $f:V\to\F_3$ is nonzero. Since all character values are
$1\pmod\varpi$ and $\sum c(\chi)=0$, one has $t\geq1$.

\begin{lemma}[Low degree after dividing by the common valuation]
\label{lem:cubic-degree}
The reduced polynomial representing $f$ on $\F_3^r$ has total degree at most
$t$. In addition,
\begin{equation}\label{eq:cubic-fparity}
 f(0)=0,\qquad f(-v)=(-1)^t f(v),\qquad\supp f\subseteq\supp H.
\end{equation}
\end{lemma}

\begin{proof}
Represent each coordinate by $0,1,2$. For $e\in\{0,1,2\}$ the polynomial
\[
 1+(\zeta^e-1)x+(\zeta^e-1)^2\frac{x(x-1)}2
\]
interpolates $x\mapsto\zeta^{ex}$ on these three representatives. Its
coefficient of $x^j$ is divisible by $\varpi^j$, for $j=0,1,2$, in
$R[1/2]$. Taking products in the coordinates and integer linear
combinations produces the reduced interpolating polynomial for $H$;
every coefficient of total degree $j$ is divisible by $\varpi^j$.

Every coefficient is also divisible by $\varpi^t$. To see this, interpolate
the values $\varpi^{-t}H(v)$ instead. The one-variable evaluation matrix
on $0,1,2$ has determinant $2$, so its inverse, and the tensor product
inverse in $r$ variables, have entries in $R[1/2]$. Thus division of the
interpolating polynomial by $\varpi^t$ leaves its coefficients in
$R[1/2]$. Reducing modulo $\varpi$ kills every monomial of degree greater
than $t$, proving the degree claim.

The support and value at zero follow from the definition. Conjugation
acts trivially on $R/(\varpi)$, while
$\overline\varpi/\varpi=-\zeta^2\equiv-1\pmod\varpi$.
Equation \eqref{eq:cubic-symmetry} therefore gives the parity identity.
\end{proof}

\begin{lemma}[Exact cost of an odd or even valuation]
\label{lem:cubic-cost}
If $t=2j+1$, then $b\geq3^j$. If $t=2j+2$, then $b\geq2\cdot3^j$.
\end{lemma}

\begin{proof}
Choose $v$ attaining the minimum defining $t$, and abbreviate $C_j(v)$ to
$C_j$. For any zero-sum integer triple the positive mass is its largest
absolute entry. Coalescing the coefficients of $c$ into this triple
cannot increase its positive mass, so
\[
 \max_j\abs{C_j}\leq b.
\]
Factor the largest common power $3^m$ from $(C_0,C_1,C_2)$, obtaining a
zero-sum integer triple $(n_0,n_1,n_2)$ not all divisible by three. Its
residue triple is either a permutation of $(0,1,2)$, or has three equal
nonzero entries.

In the first case,
$n_0+n_1\zeta+n_2\zeta^2$ is divisible by $\varpi$ but not by
$\varpi^2$, because its two coefficients in the basis $1,\zeta$ are not
both divisible by three. Thus $t=2m+1$ and
$b\geq3^m$.

In the second case put $A=n_0-n_2$, $B=n_1-n_2$. Both are divisible by
three, whereas
\[
 \frac{A+B}{3}=-n_2\not\equiv0\pmod3.
\]
Consequently $(A+B\zeta)/3$ is not divisible by $\varpi$, so $t=2m+2$.
Three integers with a common nonzero residue modulo three and zero sum
cannot all have absolute value at most one. Hence
$b\geq2\cdot3^m$.
\end{proof}

\subsection{The coding input and the minimum}

We use these established Reed--Muller facts
\cite{leducq2012,leducq2010}. On $\F_3^r$, nonzero reduced polynomials of
degree at most $2j+1<2r$ have support at least
$2\cdot3^{r-j-1}$. At degree at most $2j<2r$, the minimum is $3^{r-j}$,
and minimum supports are affine subspaces of codimension $j$. For
$1\leq j\leq r-1$, the next weight is $4\cdot3^{r-j-1}$; each support
of that weight is contained in an affine subspace of codimension $j-1$.
Minimum supports at odd degree lie in an affine subspace of codimension
$j$. The second-weight formula at degree two includes $r=2$.

\begin{lemma}[The symmetry excludes the even minimum]
\label{lem:cubic-weight}
For the function $H$ in \eqref{eq:cubic-H}, the following bounds hold:
\begin{align}
 t=2j+1<2r&\quad\Longrightarrow\quad
  3^{-r}\abs{\supp H}\geq\frac2{3^{j+1}},\label{eq:cubic-odd}\\
 t=2j<2r&\quad\Longrightarrow\quad
  3^{-r}\abs{\supp H}\geq\frac4{3^{j+1}}.
  \label{eq:cubic-even}
\end{align}
In all cases $3^{-r}\abs{\supp H}\geq2/3^r$.
\end{lemma}

\begin{proof}
For odd $t$, apply the minimum-weight bound to $f$ and use
$\supp f\subseteq\supp H$. For even $t=2j$, a minimum-weight $f$ would
have support an affine subspace $A$ of codimension $j$.
Equation \eqref{eq:cubic-fparity} gives $A=-A$. An affine subspace over
$\F_3$ invariant under negation contains zero: it contains the midpoint
of any $v$ and $-v$. This contradicts $f(0)=0$. The second-weight bound
therefore applies. Finally $\supp H$ is nonempty, invariant under
negation, and excludes zero, so it contains at least one pair
$\{v,-v\}$.
\end{proof}

\begin{proof}[Proof of Theorem~\ref{thm:cubic-min}]
The density formula is Lemma~\ref{lem:cubic-trace}. First suppose
$a\leq r-2$ and $3^a\leq k<2\cdot3^a$. Lemma~\ref{lem:cubic-cost}
forces $t\leq2a+1$. An odd value $t=2j+1$ has $j\leq a$, and
\eqref{eq:cubic-odd} gives $\delta(A,C)\geq2/3^{a+1}$.
An even value $t=2j$ has $j\leq a$ and gives the strictly larger bound
$4/3^{a+1}$.

If $a\leq r-2$ and $2\cdot3^a\leq k<3^{a+1}$, the cost lemma forces
$t\leq2a+2$. The weakest bound is now at $t=2a+2$, giving
$4/3^{a+2}$; every smaller $t$ gives a larger bound. If
$k\geq3^{r-1}$, use the universal two-point bound.

For sharpness choose a subspace $U\leq X$ of dimension $a$ and
$\theta\notin U$. The packets $U$ and $\theta+U$ have size $3^a$ and
Fourier difference
\begin{equation}\label{eq:cubic-construct-odd}
 H(v)=3^a\mathbf1_{U^\perp}(v)(1-\theta(v)).
\end{equation}
Its support has relative size $3^{-a}(2/3)$.
For the other range choose $\alpha,\beta$ independent modulo $U$ and
take the two packets
$(\alpha+U)\mathbin{\dot\cup}(-\alpha+U)$ and
$(\beta+U)\mathbin{\dot\cup}(-\beta+U)$.
For $v\in U^\perp$, write $\alpha(v)=\zeta^x$ and
$\beta(v)=\zeta^y$. Then
\begin{equation}\label{eq:cubic-construct-even}
 H(v)=3^{a+1}\bigl(\mathbf1_{x=0}-\mathbf1_{y=0}\bigr).
\end{equation}
Exactly one of $x,y$ vanishes in four of their nine equally frequent
values, giving $4/3^{a+2}$. Append identical entries to both packets to
reach the prescribed size $k$. For saturation use
\eqref{eq:cubic-construct-odd} with $\dim U=r-1$.
\end{proof}

\subsection{Rigidity of equality}

\begin{lemma}[Support in a subspace forces repeated coefficients]
\label{lem:cubic-periods}
If $H$ is supported in a linear subspace $A\leq V$, then its integral
Fourier coefficients $c$ are constant on each coset of $A^\perp\leq X$.
\end{lemma}

\begin{proof}
Fourier inversion gives
\[
 c(\chi)=3^{-r}\sum_{v\in A}H(v)\overline{\chi(v)}.
\]
Replacing $\chi$ by $\chi+u$ for $u\in A^\perp$ leaves every summand
unchanged.
\end{proof}

\begin{proof}[Proof of Theorem~\ref{thm:cubic-equality}]
In the lower nonsaturated range, equality in the preceding bounds forces
$t=2a+1$. The support of $f$ is contained in that of $H$ and already has
the minimum permitted size, so the two supports agree. The coding
classification puts this support inside an affine subspace $A\leq V$
of codimension $a$, occupying $2/3$ of it. It spans $A$, because a
proper affine subspace contains at most $1/3$ of $A$. The symmetry of
$\supp H$ implies $A=-A$, so $A$ is linear.

Let $U=A^\perp$, of dimension $a$. By
Lemma~\ref{lem:cubic-periods}, $c$ is constant on cosets of $U$. Its
positive mass is therefore a positive integer multiple of $3^a$, less
than $2\cdot3^a$. On $X/U$ its coefficient vector has one $+1$ and
one $-1$. This is the two-coset description in part~1.

In the upper nonsaturated range equality similarly forces $t=2a+2$,
and $\supp f=\supp H$ has the second Reed--Muller weight. Its support
lies in an affine subspace $A$ of codimension $a$, occupying $4/9$ of
it. Since $4/9>1/3$, this is again the affine span of the support and is
linear by symmetry. With $U=A^\perp$, the coefficient vector descends
to $X/U$. Its positive mass is at least $2\cdot3^a$ by the cost lemma
and less than $3\cdot3^a$, so the quotient consists of two disjoint
packets of size two.

If either quotient packet repeats a character, agreement of the two
traces forces all its local character values to coincide with those
of the other packet, by equality in the triangle inequality. Agreement
then lies in a proper linear subspace and has relative size at most
$1/3$. The required agreement is $5/9$, so both quotient packets have
distinct entries. Write them as $\{x,y\}$ and $\{z,w\}$, with all four
entries different. If $x+y-z-w\ne0$, agreement of the local polynomials
implies that this nontrivial character takes the value one. This again
allows agreement on at most $1/3$ of the quotient dual. Thus
$x+y=z+w$.

Division by two in $\F_3$ defines
\[
 \eta=(x+y)/2,\qquad\alpha=(x-y)/2,\qquad\beta=(z-w)/2.
\]
Distinctness of the four entries says that $\alpha,\beta$ are nonzero
and unequal up to sign. Over $\F_3$, this is their linear independence.
Lifting the quotient elements gives part~2.

At saturation, equality means $\supp H=\{v,-v\}$ for some nonzero
$v$. Let $U=\ker v<X$. Lemma~\ref{lem:cubic-periods} makes $c$
constant on the three cosets of $U$. Its three integer values sum to
zero and are not all zero. Their positive mass is their largest
absolute value, giving the formula in part~3. Conversely any such
function has Fourier support contained in $U^\perp$; its value at
zero vanishes, and nonzero conjugate values remain at the other two
points. All stated converses also follow from the explicit
constructions above.
\end{proof}

\subsection{Sharp tolerance of erased and corrupted observations}

\begin{corollary}[Exact asymptotic observation thresholds]
\label{cor:cubic-errors}
Consider products of exactly $k$ characters from a fixed rank-$r$ group
$X$, observed through their local polynomials at primes not dividing $Q$.
\begin{enumerate}
\item If an indicated erasure set has upper natural density strictly less
than $\Delta_{r,k}$, the remaining exact observations identify at most one
packet. The threshold is sharp.
\item If an arbitrary record differs from the true packet on a set of upper
natural density strictly less than $\Delta_{r,k}/2$, there is at most one
packet satisfying that error bound. This threshold is sharp as well.
\end{enumerate}
\end{corollary}

\begin{proof}
Two packets agreeing outside an erasure set would have disagreement
density bounded by that set's upper density. Two packets each within
error density $\varepsilon$ of one record would disagree on a set of
upper density at most $2\varepsilon$. Theorem~\ref{thm:cubic-min} gives
uniqueness in both cases.

For sharpness choose an extremizing pair and erase precisely its
disagreement primes. For corrupted observations, order the disagreement
primes and alternate their allocation to two sets. Each has natural
density $\Delta_{r,k}/2$. Use the first packet's value on one set and
the second packet's value on the other, and their common value elsewhere.
This single record is at the threshold distance from both packets.
\end{proof}

\subsection{Effective arithmetic consequences of Dirichlet GRH}

We now use the input Dirichlet GRH to make separation visible in a finite
prime interval. Fix a smooth nonnegative function $w$ supported in
$(1/2,1)$ with $\int_0^\infty w(u)\,du=1$. For $F:V\to\C$ set
\[
 M_x(F)=\sum_{p\nmid Q}(\log p)w(p/x)F([p]),
 \qquad \norm{F}_{\mathrm A}=\sum_{\chi\in X}\abs{\widehat F(\chi)},
\]
where $\widehat F$ uses normalized finite Fourier coefficients and
$\E_V F=3^{-r}\sum_{v\in V}F(v)$.

\begin{lemma}[Smoothed GRH interface]\label{lem:cubic-grh}
Under Dirichlet GRH, for $x\geq2$,
\begin{equation}\label{eq:cubic-grh-interface}
 \abs{M_x(F)-x\E_V F}
 \leq C_w\norm{F}_{\mathrm A}\sqrt{x}\log(2Q).
\end{equation}
The constant $C_w$ is independent of $X,Q,r,F$ and $x$.
\end{lemma}

\begin{proof}
For a character $\chi\in X$, apply the smoothed explicit formula to its
primitive inducer. The principal character contributes $x$. Under GRH
each nontrivial zero contributes a term of absolute value
$\sqrt{x}\abs{\widetilde w(1/2+i\gamma)}$, where
$\widetilde w$ is the Mellin transform. Repeated integration by parts
gives rapid decay in $\gamma$, while the usual unit-interval zero count
is $O(\log(2Q(\abs\gamma+2)))$. Their sum is
$O_w(\sqrt{x}\log(2Q))$. The trivial-zero and archimedean terms satisfy
the same bound. Subtracting prime powers of exponent at least two costs
$O_w(\sqrt{x})$, and removing primes dividing $Q$ costs
$O_w(\log Q)$. This proves
\eqref{eq:cubic-grh-interface} for each character, including imprimitive
ones. Fourier expansion and the triangle inequality prove the assertion
for $F$. This is the analytic interface used in the article's quadratic
matching result \cite{article}.
\end{proof}

\begin{theorem}[A finite distinguishing prime]\label{thm:cubic-prime}
Under Dirichlet GRH there is a constant $C_w'>0$ such that every distinct
pair of size-$k$ packets from $X$ has a prime with different local
polynomials in $(x/2,x)$ whenever
\[
 x\geq C_w' k^2\log^2(2Q).
\]
\end{theorem}

\begin{proof}
Let $B=\sum_\chi\abs{c(\chi)}=2b\leq2k$ and use $F=\abs H^2$.
Parseval and integrality give
\[
 \E_V\abs H^2=\sum_\chi c(\chi)^2\geq B,\qquad
 \norm{\abs H^2}_{\mathrm A}\leq B^2.
\]
Hence $M_x(\abs H^2)\geq Bx-C_wB^2\sqrt{x}\log(2Q)>0$ once the
constant $C_w'$ is large enough. A positive summand supplies the
required prime, and Lemma~\ref{lem:cubic-trace} identifies it as a
local-polynomial disagreement.
\end{proof}

The $k^2\log^2 Q$ scale follows from the general Fourier argument. The
complete extremizer classification gives a stronger uniform statement
for counting the distinguishing primes of an extremal pair.

\begin{theorem}[Exact norms and uniform counting for all extremizers]
\label{thm:cubic-norm}
Let $S\subset V$ be the disagreement set of a pair attaining
$\Delta_{r,k}$. Then
\begin{equation}\label{eq:cubic-exactnorm}
 \norm{\mathbf1_S}_{\mathrm A}=
 \begin{cases}
  4/3,&\text{the lower nonsaturated range or the saturated range},\\
  16/9,&\text{the upper nonsaturated range}.
 \end{cases}
\end{equation}
Consequently, under Dirichlet GRH,
\begin{equation}\label{eq:cubic-extcount}
 \sum_{\substack{p\nmid Q\\P_A(p,T)\ne P_C(p,T)}}
    (\log p)w(p/x)
 =\Delta_{r,k}x+O_w\bigl(\sqrt{x}\log(2Q)\bigr),
\end{equation}
with a constant independent of packet size and rank.
\end{theorem}

\begin{proof}
The common translation $\eta$ multiplies $H$ by a unit-modulus character
and does not change its support. In the lower range, put
$K=U+\langle\theta\rangle$. Then
\[
 \mathbf1_S=\mathbf1_{U^\perp}-\mathbf1_{K^\perp}.
\]
The normalized Fourier transform of $\mathbf1_{L^\perp}$ is
$\abs{L}^{-1}\mathbf1_L$. Thus the coefficients above have magnitude
$2/3^{a+1}$ on $U$, magnitude $1/3^{a+1}$ on the other two cosets of
$U$ in $K$, and vanish elsewhere. Their absolute sum is $4/3$.

For the upper range let
\[
 K_1=U+\langle\alpha\rangle,\quad
 K_2=U+\langle\beta\rangle,\quad
 K=U+\langle\alpha,\beta\rangle.
\]
Equation \eqref{eq:cubic-construct-even} gives
\[
 \mathbf1_S=\mathbf1_{K_1^\perp}+\mathbf1_{K_2^\perp}
             -2\mathbf1_{K^\perp}.
\]
On the nine cosets of $U$ in $K$, the absolute Fourier coefficients
are $4/3^{a+2}$ on $U$, $1/3^{a+2}$ on four cosets, and
$2/3^{a+2}$ on the remaining four. Their sum is
$(4+4+8)/9=16/9$.

At saturation $S=U^\perp\setminus\{0\}$ with $\dim U=r-1$;
the first norm calculation applies with $K=X$. Finally apply
Lemma~\ref{lem:cubic-grh} to $\mathbf1_S$.
\end{proof}

Equation \eqref{eq:cubic-extcount} supplies a finite observation guarantee:
after division by $x$, the weighted disagreement estimate differs from
the exact density by $O_w(x^{-1/2}\log(2Q))$, uniformly over every
extremizer. The algebraic classification is responsible for this
uniformity.
